% Tree of Knowledge: Specialization Depth, Adapter Rank, and Storage Tier
% Are the Same Dimension
%
% The central claim: in a well-designed MoE architecture, how specialized
% an expert is (information-theoretic), how large its adapter is (parameters),
% and where it lives in the memory hierarchy (hardware) are not three separate
% design decisions — they are one.

\documentclass{article}
\usepackage[preprint]{neurips_2024}
\usepackage[utf8]{inputenc}
\usepackage[T1]{fontenc}
\usepackage{hyperref}
\usepackage{url}
\usepackage{booktabs}
\usepackage{amsfonts}
\usepackage{amsmath}
\usepackage{nicefrac}
\usepackage{microtype}
\usepackage{graphicx}
\usepackage{xcolor}

\newcommand{\placeholder}[1]{\textcolor{red}{\textbf{[#1]}}}
\newcommand{\tok}{Tree of Knowledge}

\title{Tree of Knowledge: Differentiating Fluid Intelligence \\ from Modular Crystallized Knowledge in Language Models}

\author{
  Erik de Bruijn \\
  Independent Researcher \\
  \texttt{erik@erikdebruijn.nl}
}

\begin{document}
\maketitle

% ============================================================================
\begin{abstract}
Current language models entangle \emph{fluid intelligence}---general reasoning and language understanding---with \emph{crystallized intelligence}---domain-specific knowledge accumulated during training. We propose the \tok{} framework, which separates these into a shared dense core and a hierarchy of progressively specialized LoRA adapters that branch like a tree through the model's depth. The central insight is that three apparently separate quantities---\textbf{specialization depth} (how different an expert's knowledge is from the generalist core), \textbf{adapter rank} (how many parameters the expert needs), and \textbf{storage tier} (where the expert lives in the memory hierarchy)---are the same dimension. A rank-1 adapter is a lens: a 4KB correction to the generalist, always cached. A rank-64 adapter is a crystal: a 256KB specialist that reshapes computation fundamentally, loaded from flash on demand. This unification arises naturally when experts are created by \emph{learntropy-driven splitting}: a teacher model monitors each expert's prediction difficulty distribution, and when it detects bimodal learntropy---some tokens easy, others hard---the expert splits. Each child inherits the parent's weights plus a LoRA perturbation, starting at rank-1 and growing as the data demands. A contrastive loss accelerates divergence between siblings. The tree grows through the model's depth: early layers share a common trunk, middle layers branch into broad domains, and deep layers fork into narrow specialties. We validate on Qwen3-1.7B (28 layers, 2.0B parameters) and ground our proposal in systematic negative results from six prior experimental arms that establish what does and does not produce genuine expert differentiation.
\end{abstract}

% ============================================================================
\section{Introduction}

The promise of Mixture-of-Experts is that different experts learn different things. We spent six experimental arms and 40+ hours of GPU training discovering that this promise is remarkably hard to fulfill. Copy-initialized experts trained on disjoint data remain virtually identical (cosine similarity 0.998 after 12M tokens). Prescriptive cost losses shift routing but the effect dilutes. Biological mechanisms help marginally. Without explicit cost signals, no specialization emerges at all.

These failures taught us something important: \textbf{expert differentiation requires three forces acting together}:
\begin{enumerate}
    \item \textbf{Different starting points.} Copy-initialized experts share an identical basin of attraction that data alone cannot escape. LoRA adapters, randomly initialized, start from genuinely different positions in parameter space.
    \item \textbf{Different training data.} Disjoint data partitions force each expert to encounter different input distributions. Without this, all experts converge to the same function.
    \item \textbf{Explicit divergence pressure.} Even with different data, convergence toward the shared optimum is the default. A contrastive loss that penalizes similarity between sibling experts provides the gradient signal that accelerates differentiation.
\end{enumerate}

The \tok{} framework instantiates all three through a single mechanism: \emph{learntropy-driven progressive splitting with variable-rank LoRA adapters}.

\subsection{The Unifying Insight}

We observe that three design decisions typically treated independently---specialization depth, adapter rank, and storage tier---are manifestations of the same underlying quantity: \emph{how different this expert's computation is from the generalist core}.

\begin{table}[h]
\centering
\caption{Specialization depth, adapter rank, and storage tier are the same dimension.}
\label{tab:unification}
\begin{tabular}{@{}llll@{}}
\toprule
Specialization & Adapter & Size & Storage \\
\midrule
Minimal (lens) & Rank-1 LoRA & 4\,KB & CPU cache / always resident \\
Moderate & Rank-16 LoRA & 64\,KB & VRAM cache \\
Substantial & Rank-64 LoRA & 256\,KB & flash (instant load) \\
Fundamental (crystal) & Full expert & 12\,MB & flash (prefetchable) \\
\bottomrule
\end{tabular}
\end{table}

A rank-1 adapter is a \emph{lens}: it makes a small correction to the generalist's output, like reading glasses that sharpen an already-clear image. A rank-64 adapter is a new \emph{crystal}: it computes a fundamentally different function, like a prism that decomposes white light into a spectrum the generalist cannot see. The rank is not a hyperparameter to tune---it is a \emph{measurement} of how different this knowledge is from the core.

% ============================================================================
\section{Architecture: Progressive LoRA Forking}

\subsection{Shared Trunk, Branching Depth}

A standard transformer with $L$ layers processes every token through the same $L$ feed-forward networks. In \tok{}, early layers remain shared (the ``trunk''), while later layers progressively branch:

\begin{itemize}
    \item \textbf{Layers 1--$k_1$}: Shared FFN. All tokens follow the same path. This is the fluid intelligence core---general reasoning, syntax, attention patterns. Always resident in VRAM.
    \item \textbf{Layers $k_1$--$k_2$}: First fork. 2--4 broad domain branches, each implemented as a LoRA adapter on the shared FFN. The router selects which branch based on the hidden state.
    \item \textbf{Layers $k_2$--$L$}: Progressive forking. Each branch may split further as specialization demands. The tree grows deeper where the data requires it.
\end{itemize}

For Qwen3-1.7B (28 layers), a natural partition is:
\begin{itemize}
    \item Layers 1--14: trunk (shared, 1.0B parameters)
    \item Layers 15--21: first fork, 2--4 branches ($\sim$2\,MB LoRA per branch)
    \item Layers 22--28: deep forks, up to 8+ branches ($\sim$0.5\,MB LoRA per branch)
\end{itemize}

Total specialist parameters: $\sim$6\,MB for 8 specialists, versus 800+\,MB for 8 full expert copies in standard MoE.

\subsection{Learntropy-Driven Splitting}

An expert splits when a teacher model detects \textbf{bimodal learntropy}: the expert handles some tokens easily (low prediction error) and others poorly (high prediction error). This bimodality is evidence that one expert is serving two distinct populations that would benefit from separate treatment.

\begin{enumerate}
    \item The teacher scores each expert's tokens by per-token loss.
    \item If the loss distribution is bimodal (measured by bimodality coefficient or Hartigan's dip test), the expert is a split candidate.
    \item The expert's LoRA adapter is duplicated. Each child inherits the parent's weights plus a small random perturbation.
    \item Children start at rank-1. If the data requires more capacity, rank grows adaptively.
    \item A contrastive loss between siblings accelerates divergence.
\end{enumerate}

The split is not on a timer or a fixed schedule---it is triggered by an information-theoretic signal from the data itself. This is Piaget's accommodation mechanism made precise: new cognitive structure is created when prediction failure reveals that existing structure is inadequate.

\subsection{Variable Rank as Specialization Measure}

After splitting, each child adapter begins at rank-1 (minimal deviation from parent). Training proceeds, and the rank grows when the reconstruction error of the low-rank adapter exceeds a threshold---the adapter needs more capacity to express the specialization the data demands.

\begin{itemize}
    \item \textbf{Rank stays at 1--4}: The specialization is a small correction. The knowledge is ``close to'' the parent. The adapter is a lens.
    \item \textbf{Rank grows to 16--64}: The specialization is substantial. The expert handles fundamentally different computation. The adapter is a crystal.
    \item \textbf{Rank approaches full}: The expert has diverged so far from the parent that it should be promoted to a full expert (or the parent should be restructured).
\end{itemize}

The rank IS the measurement of specialization depth. It is not chosen by a human---it emerges from the data.

\subsection{Contrastive Divergence Loss}

Without explicit pressure, sibling experts drift toward the shared optimum (we observed CosSim 0.998 after 12M tokens of disjoint data training). The contrastive loss provides the missing force:

\begin{equation}
\mathcal{L}_{\text{contrast}} = \lambda_c \sum_{i < j} \max(0, \cos(\theta_{ij}) - \tau)
\end{equation}

where $\theta_{ij}$ is the angle between sibling expert weight vectors and $\tau$ is a target maximum similarity. Combined with disjoint training data, this produces rapid differentiation.

% ============================================================================
\section{Why These Three Are the Same Dimension}

The unification is not a convenient mapping---it is a structural identity:

\textbf{Information-theoretic view.} An expert that handles tokens from a narrow, homogeneous distribution needs few additional parameters to deviate from the generalist (low KL divergence from parent $\rightarrow$ low rank). An expert that handles a fundamentally different distribution needs many parameters (high KL divergence $\rightarrow$ high rank).

\textbf{Hardware view.} Parameters that are rarely needed should live on slower, cheaper storage. Rank-1 adapters (4\,KB) are so small they never leave the cache. Rank-64 adapters (256\,KB) load in microseconds from flash. The rank determines the storage cost because it determines the size, and the size determines the tier.

\textbf{Frequency view.} Common patterns are handled by the shared trunk (always resident). Moderately common domains are handled by low-rank adapters (VRAM-cached). Rare specializations are handled by high-rank adapters (flash-loaded on demand). The Zipf distribution of language ensures that most tokens need no specialist at all.

These three views produce the same ordering: common $\rightarrow$ rare = low rank $\rightarrow$ high rank = always cached $\rightarrow$ loaded on demand. \textbf{It is one dimension, not three.}

% ============================================================================
\section{Lessons from Six Experimental Arms}
\label{sec:lessons}

\begin{table}[h]
\centering
\caption{Six arms that taught us what doesn't work, motivating every design decision in \tok{}.}
\label{tab:lessons}
\begin{tabular}{@{}llp{6cm}@{}}
\toprule
Arm & Result & Lesson for \tok{} \\
\midrule
A: Prescriptive cost & Gini 0.60$\to$0.49 & Cost loss creates structure but it dilutes \\
B: Reuse-distance proxy & Gini 0.05 & Descriptive loss fails entirely \\
C: Adaptive + cost & 25 rounds, Gini dilutes & Piaget mechanism works, but routing skew is wrong metric \\
D: Pure data-driven & Gini 0.11 & No specialization without resource constraint \\
E: Bio-integrated & Gini 0.46, marginal & Six mechanisms help modestly, not fundamentally \\
KD-Warm upcycle & CosSim 0.998 & Copy-init experts don't differentiate via data alone \\
\bottomrule
\end{tabular}
\end{table}

The KD-Warm experiment is the most decisive: after 12M tokens of forced-disjoint training on Qwen3-1.7B, expert weight cosine similarity remained at 0.998. The copy initialization creates a basin of attraction that disjoint data cannot escape on practical timescales. This motivates LoRA adapters (different starting points) and contrastive loss (explicit divergence pressure).

% ============================================================================
\section{Experimental Plan}

We validate on Qwen3-1.7B (28 layers, 2.0B parameters, C4 data):

\begin{enumerate}
    \item \textbf{Baseline.} Original dense model PPL on C4 validation set.
    \item \textbf{Trunk identification.} Freeze layers 1--14 (trunk), add rank-1 LoRA adapters to layers 15--28, train on full data. Measure: does quality recover? Which layers need adaptation?
    \item \textbf{First split.} Using learntropy bimodality, identify the first split candidate. Create two sibling LoRA adapters with contrastive loss. Measure: CosSim between siblings over training (target: $<$0.95 within 1000 steps).
    \item \textbf{Progressive forking.} Allow the tree to grow through multiple splits. Measure: how many specialists emerge? What ranks do they reach? Does the tree structure correspond to interpretable domains?
    \item \textbf{Variable rank validation.} Start all adapters at rank-1. Allow rank growth when reconstruction error exceeds threshold. Measure: final rank distribution across experts. Hypothesis: follows a power law (few high-rank specialists, many low-rank lenses).
    \item \textbf{Hot-loading.} Remove a specialist, verify quality drops only on its domain. Add it back, verify recovery. This validates modularity.
\end{enumerate}

% ============================================================================
\section{Related Work}

\textbf{LoRA and adapters.} Hu et al.~\cite{hu2022lora} introduced Low-Rank Adaptation for efficient fine-tuning. We extend LoRA from a fine-tuning tool to a \emph{structural unit of specialization} within a single model.

\textbf{MoE upcycling.} Converting dense models to MoE by copying FFN weights~\cite{komatsuzaki2023sparse}. Our KD-Warm experiment shows this produces identical experts. Drop-Upcycling~\cite{dropupcycle2025} addresses this with random re-initialization; we propose LoRA + contrastive loss as a principled alternative.

\textbf{Geometric routing.} Hash Layers~\cite{roller2021hash}, EMoE~\cite{emoe2026}, IDA-MoE~\cite{idamoe2025} partition the embedding space for routing. Our analysis shows raw hidden states lack geometric structure; the routing must be learned, with the tree providing initialization.

\textbf{Specialization theory.} B\'{e}na \& Goodman~\cite{bena2025modularity} prove resource constraints are necessary for specialization. Our contrastive loss provides the constraint; the rank measures the result.

\textbf{Low-rank routing.} L2R~\cite{l2r2026} projects to $r{=}2$ dimensions for routing, confirming that routing decisions are intrinsically low-dimensional. Our variable-rank adapters extend this: the routing is low-rank, and so is the specialization itself.

% ============================================================================
\section{Conclusion}

Six failed experiments, one working hypothesis: specialization depth, adapter rank, and storage tier are the same dimension. The \tok{} framework exploits this identity through progressive LoRA forking driven by learntropy splitting signals. Each specialist is born as a rank-1 lens and grows into whatever the data demands---a rank-4 correction, a rank-64 crystal, or a full expert. The tree structure is not imposed; it emerges from the information structure of language, growing where knowledge is dense and diverse, staying shallow where a generalist suffices.

The practical consequence is a model architecture where specialists are measured in kilobytes, loadable in microseconds, trainable independently on commodity hardware, and composable by a community of contributors. Not because we optimized for these properties, but because they are the natural expression of a single underlying principle: \emph{the more specialized the knowledge, the more parameters it needs, and the rarer it is needed}.

% ============================================================================
\bibliographystyle{plainnat}
\begin{thebibliography}{99}
\bibitem{hu2022lora} E. Hu et al. LoRA: Low-Rank Adaptation of Large Language Models. \textit{ICLR}, 2022.
\bibitem{komatsuzaki2023sparse} A. Komatsuzaki et al. Sparse Upcycling. \textit{arXiv:2212.05055}, 2023.
\bibitem{dropupcycle2025} Drop-Upcycling: Training Sparse MoE with Partial Re-initialization. \textit{ICLR}, 2025.
\bibitem{roller2021hash} S. Roller et al. Hash Layers For Large Sparse Models. \textit{NeurIPS}, 2021.
\bibitem{emoe2026} EMoE: Eigenbasis-Guided Routing. \textit{arXiv:2601.12137}, 2026.
\bibitem{idamoe2025} IDA-MoE: Decoupled Routing via GMMs. \textit{ACM MM}, 2025.
\bibitem{bena2025modularity} G. B\'{e}na, D. Goodman. Dynamics of Specialization Under Resource Constraints. \textit{Nature Communications}, 2025.
\bibitem{l2r2026} L2R: Low-Rank Lipschitz-Controlled Routing. \textit{arXiv:2601.21349}, 2026.
\bibitem{piaget1952} J. Piaget. \textit{The Origins of Intelligence in Children}. 1952.
\bibitem{vygotsky1978} L. Vygotsky. \textit{Mind in Society}. Harvard University Press, 1978.
\end{thebibliography}

\end{document}
